\documentclass[10pt,a4paper]{amsart}
\usepackage[margin=19mm]{geometry}
\usepackage{amsmath,amssymb,mathtools}
\usepackage[scr=rsfs,cal=euler]{mathalfa}
\usepackage{xfrac}
\usepackage[unicode,hidelinks,bookmarks=false]{hyperref}
\newcommand{\ie}{i.e.\ }
\newcommand{\cf}{cf.\ }
\newcommand{\resp}{respectively\ }
\newcommand{\Subsection}{\subsection}
\providecommand{\nobreakdash}{\nobreak\hbox{-}}
\setlength{\emergencystretch}{2em}
\multlinegap0pt
\usepackage{bm}
\usepackage{bbm}
\usepackage{dsfont}
\usepackage{xspace}
\usepackage{cancel}
\usepackage{mathtools}
\usepackage{amsthm}
\makeatletter
\@ifundefined{theorem}{\newtheorem{theorem}{Theorem}}{}
\@ifundefined{lemma}{\newtheorem{lemma}[theorem]{Lemma}}{}
\@ifundefined{proposition}{\newtheorem{proposition}[theorem]{Proposition}}{}
\makeatother
\newcommand\NN{\mathbb{N}}
\newcommand\RR{\mathbb{R}}
\newcommand\diam{\mathrm{diam}}
\newcommand{\N}{\mathbb{N}}
\newcommand{\R}{\mathbb{R}}
\title[Equality of H\"older exponents for distribution functions of]{Equality of H\"older exponents for distribution functions of Gibbs measures}
\author{Pieter Allaart, Johannes Jaerisch}
\date{Source: arXiv:2509.11527v3 (CC BY 4.0)}
\begin{document}
\maketitle

\noindent\emph{Source and licence.} Equality of H\"older exponents for distribution functions of Gibbs measures. Version arXiv:2509.11527v3 (CC BY 4.0), distributed under the Creative Commons Attribution 4.0 International licence, as recorded in the arXiv licence metadata for this exact version and shown on the abstract page https://arxiv.org/abs/2509.11527v3. Full rights notice: licenses/arxiv-2509.11527-CC-BY-4.0.txt.\par\smallskip

\noindent\textit{Editorial scope.} Counted unit: the Proposition whose source label is \texttt{prop:no-positive-derivative} in the article (source0.tex lines 472--478, source label \texttt{prop:no-positive-derivative}), delivered together with the article's own complete original proof of it (source0.tex lines 480--633, one \texttt{proof} environment of 154 lines). No mathematics is written, repaired or re-derived: statement and proof are the article's own text line for line. Required context retained uncounted: eq:bdd (lines 288--290); eq:gibbs (lines 291--293); lem:non-adjacency (lines 430--432); lem:cohom-1 (lines 446--461). Every definition, display and label of those lines is retained unchanged. Omitted from this package and not redistributed: 1--287, 294--429, 433--445, 462--471, 479--479, 634--1036. Nothing inside a retained excerpt is omitted or altered: every retained body is delivered exactly as the article's lines are written, so no omission receipt of any kind accompanies this package. Citations: the retained excerpts carry no citation command at all, so no bibliography entry of the article is transcribed and no reference is redistributed. Reference adaptation: none was needed and none was made; every reference inside the delivered unit resolves inside it, so no unresolved reference or citation is produced. Printed slips kept literal: no printed word, symbol, hypothesis or number of the counted statement or of its proof was altered. Layout and class: the article's own class is \texttt{amsart} with packages including amsmath, amssymb, enumerate, mathrsfs, amsthm, xcolor; the wrapper is the lane's pinned \texttt{amsart} editorial wrapper at 10pt on A4, which restates the theorem-like environments the retained text needs and adds the article's own macro definitions that the retained text needs (\texttt{\textbackslash N}, \texttt{\textbackslash NN}, \texttt{\textbackslash R}, \texttt{\textbackslash RR}, \texttt{\textbackslash diam}), with extra wrapper packages (bm, bbm, dsfont, xspace, cancel, mathtools). The delivered numbering is the unit's own. Assets: no retained span contains a figure environment, a graphics-inclusion command, a PDF-page inclusion, a table or a TikZ picture; no image file is copied into this package, the package declares no asset dependency and no image file is redistributed with it. Licence: version arXiv:2509.11527v3 of this article is distributed under the Creative Commons Attribution 4.0 International licence, as recorded in the arXiv licence metadata for this exact version and shown on the abstract page https://arxiv.org/abs/2509.11527v3. A full rights notice is delivered at licenses/arxiv-2509.11527-CC-BY-4.0.txt. Characters: every retained excerpt is pure ASCII, so the delivered engine, XeLaTeX with its default Latin Modern text font, reports no missing character.
\begin{equation}\label{eq:bdd}
\sup_{n\in \N} \sup_{w\in I^n} \sup_{\rho,\tau\in I^\N }\left|S_n \psi(w\rho)-S_n\psi(w\tau)\right|<\infty.
\end{equation} We say that a Borel probability measure $\mu_\psi$ on $I^{\N}$  is a Gibbs measure for $\psi$  if there exist two constants $C\ge 1$ and $P\in \R$ such that for each $n\in \N$ and for all $w \in I^n$ we have 

\begin{equation}\label{eq:gibbs}
C^{-1}\le \frac{\mu_\psi[w_1\dots w_n]}{e^{S_n\psi(\overline w)-nP }}\le C.
\end{equation} Here $\overline{w}$ denotes the infinite periodic sequence with period block $w$. 

\begin{lemma} \label{lem:non-adjacency}
Let {$\ell\geq 3$} and let $\tau_1\dots\tau_\ell\in I^\ell$ be a word such that $\tau_1\neq\tau_2$. Then for any $\omega\in I^\NN$ there are infinitely many integers $n$ such that the cylinder intervals $\pi([\omega|_{n+3\ell}])$ and $\pi([\omega|_{n}(\tau_1\dots\tau_\ell)^3])$ do not have a common endpoint.
\end{lemma}

\begin{lemma} \label{lem:cohom-1} 
Suppose there exists $\ell_0\ge 1$ such that the set
\[
A:=\left\{\frac{S_{\ell}\psi\left(\overline{\tau_{1}\dots \tau_{\ell}}\right)}{S_{\ell}\varphi\left(\overline{\tau_{1}\dots \tau_{\ell}}\right)}:\quad \ell \ge \ell_0,\quad \tau_1, \dots, \tau_\ell \in I \right\}
\]
is finite.
%\textcolor{red}{(remove) Suppose that there
%exists a finite set $A\subset\mathbb{R}$ and $\ell_{0}\ge1$ such
%that for all $\ell\ge\ell_{0}$ and all {$\tau_{1},\dots,\tau_{\ell}\in I$}
%we have 
%\[
%\frac{S_{\ell}\psi\left(\overline{\tau_{1}\dots \tau_{\ell}}\right)}%{S_{\ell}\varphi\left(\overline{\tau_{1}\dots \tau_{\ell}}\right)}\in A.
%\]
%}
Then $\psi$ is cohomologous to $(\dim_{H}\Lambda)\varphi$  and $ A=\left\{\dim_{H}\Lambda\right\} $.
\end{lemma}

\begin{proposition} \label{prop:no-positive-derivative}
For every $k\in\mathbb{N}$ and $x\in\mathbb{R}$ the limit 
\[
\lim_{y\rightarrow x}\frac{F(y)-F(x)}{(y-x)^{k}}
\]
either does not exist, or it takes the value $0$ or $+\infty$. 
\end{proposition}

\begin{proof}%[Proof of Proposition \ref{prop:no-positive-derivative}]
{If $x\in\RR\backslash\Lambda$, then $F$ is constant on a neighborhood of $x$ since $\Lambda$ is closed, and hence the limit in the proposition is equal to $0$. Therefore, we may assume that $x\in\Lambda$.}
The statement is trivial for even $k$, since $F$ is increasing. So below we will assume that $k$ is odd.

Suppose by way of contradiction that 
\[
\lim_{y\rightarrow x}\frac{F(y)-F(x)}{(y-x)^{k}}=L\in(0,+\infty).
\]
Note that the limit cannot be negative, since $F$ is increasing.

Let $\omega \in I^\N$ such that $\pi(\omega)=x$.
%If $s_{n}\le x\le t_{n}$ then 
%\[
%0\le\frac{t_{n}-x}{t_{n}-s_{n}}\le1.
%\]
Let $[s_n,t_n]:=\pi\left(\left[\omega|_{n}\right]\right)$. Then $s_n\leq x\leq t_n$, and
\begin{align}
\begin{split}
\frac{F(t_{n})-F(s_{n})}{(t_{n}-s_{n})^{k}} & =\left(\frac{t_{n}-x}{t_{n}-s_{n}}\right)^{k}\cdot\frac{F(t_{n})-F(x)}{(t_{n}-x)^{k}}\\
 & \quad+\left(\frac{x-s_{n}}{t_{n}-s_{n}}\right)^{k}\cdot\frac{F(s_{n})-F(x)}{(s_{n}-x)^{k}}.
\end{split}
\label{eq:secant-slope}
\end{align}
Denote 
$$r_{n}:=\frac{t_{n}-x}{t_{n}-s_{n}}.$$ 
Observe that $r_n\in[0,1]$ and
$$1-r_{n}=\frac{x-s_{n}}{t_{n}-s_{n}}.$$

There exists  $\ell\in\mathbb{N}$ and $\tau_{1}\dots \tau_{\ell}\in I^\ell$
containing at least two different letters such that 
\[
S_{\ell}(\psi-k\varphi)\left(\overline{\tau_{1}\dots \tau_{\ell}}\right)\neq0.
\]
Indeed such $\tau_{1}\dots \tau_{\ell}$ always exist for otherwise, with the set  $A$ in Lemma \ref{lem:cohom-1} given by the finite set $\{k\}\cup\left\{ \psi(\overline{i})/\varphi(\overline{i}):i\in I\right\}$, Lemma \ref{lem:cohom-1} implies that $\psi$ is cohomologous to $(\dim_{H}\Lambda)\varphi$
and $\dim_{H}\Lambda=k=1$ {(since $\Lambda\subset\RR$)}, contradicting our hypothesis. Replacing $\tau_1\dots\tau_\ell$ by a cyclical permutation if necessary, we may assume that $\tau_2\neq\tau_1$. {Further, we may assume without loss of generality that $\ell\geq 3$: If $\ell=2$, we can replace $\ell$ and $\tau_1\tau_2$ with $\ell'=4$ and $\tau_1'\dots\tau_4'=(\tau_1\tau_2)^2$, and achieve the same effect.}

%Since $\tau_{1}\dots \tau_{\ell}$ contains two different letters, there exist infinitely many $n\ge1$ such that 
%\begin{equation} \label{eq:different-words}
%\omega_{n+1}\dots\omega_{n+\ell}\neq \tau_{1}\dots \tau_{\ell}.
%\end{equation}
%From now, we restrict our attention to such $n$. 
Taking a subsequence if necessary, we may assume that $\lim r_{n}=r\in[0,1]$.
Then by \eqref{eq:secant-slope}, we have
\[
\frac{F(t_{n})-F(s_{n})}{(t_{n}-s_{n})^{k}}\rightarrow\left(r^{k}+(1-r)^{k}\right)L\in\left[2^{1-k}L,L\right],
\]
as routine calculus shows that the function $r\mapsto r^k+(1-r)^k$ is minimized at $r=1/2$.
Hence,
%there exist $c_{1},c_{2}>0$ (depending only on $k$) such that \textcolor{cyan}{(copy second displayed formula of page 7 here?, $c_1, c_2$ are explicit ?)}
\begin{equation} \label{eq:bounded-ratio}
\lim_{n\rightarrow\infty}\frac{\mu_\psi\left(\left[\omega|_{n}\right]\right)}{\diam\left(\pi\left(\left[\omega|_{n}\right]\right)\right)^{k}}=\lim_{n\rightarrow\infty}\frac{F(t_{n})-F(s_{n})}{(t_{n}-s_{n})^{k}}\in\left[2^{1-k}L,L\right]. %\label{eq:bdd}
\end{equation}

Now, we fix an integer {$m\geq 3$} and consider the interval
\[
\pi\left(\left[\omega|_{n}(\tau_{1}\dots \tau_{\ell})^{m}\right]\right)=:\left[s_{n,m},t_{n,m}\right].
\]
Put
\[
r_{n,m}:=\frac{t_{n,m}-x}{t_{n,m}-s_{n,m}}.
\]
By Lemma \ref{lem:non-adjacency} there is an infinite subset $\mathcal{N}\subset\NN$ such that for each $n\in\mathcal{N}$, the point $x$ and the interval $[s_{n,m},t_{n,m}]$ are separated by a full cylinder interval of level $n+3\ell$. 
Replacing $\mathcal{N}$ by an infinite subset if necessary, we may assume that either $t_{n,m}>x$ for all $n\in \mathcal{N}$, or $t_{n,m}<x$ for all $n\in \mathcal{N}$. Without loss of generality, we assume the former. The argument in the latter case is similar. Thus, $r_{n,m}>0$ for all $n\in\mathcal{N}$.

Let $r_{\min}>0$ be a number such that $|\phi_i'(x)|\geq r_{\min}$ for all $i\in I$ and $x\in X$. Then
\begin{equation} \label{eq:diameters}
\diam (\pi([uv]))\geq \diam (\pi([u]))\cdot r_{\min}^{|v|}
\end{equation}
for any two finite words $u$ and $v$.

For $n\in\mathcal{N}$, let $[u_n,v_n]$ be a cylinder interval of length $n+3\ell$ separating $x$ and $[s_{n,m},t_{n,m}]$. Then by \eqref{eq:diameters},
\[
|t_{n,m}-x|\geq v_n-u_n\geq \diam (\pi([\omega|_n]))\cdot r_{\min}^{3\ell}.
\]
Hence, there is a constant $C=C(\ell)>1$, depending only on $\ell$ but not on $n$ or $m$, such that
\[
C^{-1}<\frac{t_{n,m}-x}{\diam\big(\pi\left(\left[\omega|_{n}\right]\right)\big)}<C,
\]
for all $n\in\mathcal{N}$. We denote this by
\begin{equation} \label{eq:comparable}
t_{n,m}-x\asymp_\ell\diam\big(\pi\left(\left[\omega|_{n}\right]\right)\big). 
\end{equation}

Now \eqref{eq:comparable} implies 
\begin{equation}
\label{eq:rn-bdd}
r_{n,m}=\frac{t_{n,m}-x}{t_{n,m}-s_{n,m}}\asymp_\ell\frac{\diam\left(\pi\left(\left[\omega|_{n}\right]\right)\right)}{\diam\big(\pi\left(\left[\omega|_{n}(\tau_{1}\dots \tau_{\ell})^{m}\right]\right)\big)}\asymp_\ell e^{-mS_{\ell}\left(\varphi\right)(\overline{\tau_{1}\dots \tau_{\ell}})}.
\end{equation}
%where the implied constants are independent of $n$ and $N$.
Hence, by passing to a subsequence, using a diagonal argument, we may assume that $r_{n,m}$ converges to some number $q_m \in \mathbb{R}$ for every $m\ge 1$.  {Since $r_{n,m}>0$ for every $n\in\mathcal{N}$, it follows that $q_m\geq 0$.}

Using {the analogue of \eqref{eq:secant-slope} with $s_{n,m}$ and $t_{n,m}$ in place of $s_n$ and $t_n$} and letting $n\rightarrow \infty$ we obtain 
\begin{equation} \label{eq:limit-bounded-below}
\lim_{n\to\infty}\frac{F(t_{n,m})-F(s_{n,m})}{(t_{n,m}-s_{n,m})^k} =\left(q_{m}^k+(1-q_{m})^k\right)L \ge 2^{1-k}L.
\end{equation}

On the other hand, by the Gibbs property \eqref{eq:gibbs} of $\mu_\psi$, the bounded distortion property \eqref{eq:bdd} and \eqref{eq:bounded-ratio} we have  
\begin{align*}
\frac{F(t_{n,m})-F(s_{n,m})}{(t_{n,m}-s_{n,m})^k} & =\frac{\mu_\psi\left(\left[\omega|_{n}(\tau_{1}\dots \tau_{\ell})^{m}\right]\right)}{\left(\diam(\pi\left(\left[\omega|_{n}(\tau_{1}\dots \tau_{\ell})^{m}\right])\right)\right)^k}\\
 & \asymp_\ell\frac{\mu_\psi\left(\left[\omega|_{n}\right]\right)}{\left(\diam(\pi\left(\left[\omega|_{n}\right])\right)\right)^k}\cdot e^{mS_{\ell}\left(\psi-k\varphi\right)(\overline{\tau_{1}\dots \tau_{\ell}})}\\
 & \asymp_\ell e^{mS_{\ell}\left(\psi-k\varphi\right)(\overline{\tau_{1}\dots \tau_{\ell}})}.
\end{align*}
%where the implied constants are independent of $n$ and $m$. 
If $k=1$ this already gives a contradiction, since we then have equality in \eqref{eq:limit-bounded-below}, which implies $S_{\ell}\left(\psi-k\varphi\right)(\overline{\tau_{1}\dots \tau_{\ell}})=0$.

Suppose now that $k\geq 3$. We have shown that
\[
2^{1-k}L\le \big(q_{m}^k+(1-q_{m})^k\big)L= \lim_{n\to\infty}\frac{F(t_{n,m})-F(s_{n,m})}{(t_{n,m}-s_{n,m})^k} \asymp_{\ell} e^{mS_{\ell}\left(\psi-k\varphi\right)(\overline{\tau_{1}\dots \tau_{\ell}})}.
\]
Since $S_{\ell}(\psi-k\varphi)\left(\overline{\tau_{1}\dots \tau_{\ell}}\right)\neq0$
it follows that $S_{\ell}(\psi-k\varphi)\left(\overline{\tau_{1}\dots \tau_{\ell}}\right)>0$, and hence, $q_{m}\rightarrow\infty$ as $m\rightarrow\infty$.
Since $k$ is odd, $r^k+(1-r)^k=kr^{k-1}+O(r^{k-2})$ as $r\to\infty$, so we conclude that $q_m^k+(1-q_m)^k\asymp q_m^{k-1}$. Thus,
%This implies $1-3q_{N}+3(q_{N})^{2}\asymp(q_{N})^{2}$. So,
\[
q_{m}^{k-1}\asymp e^{mS_{\ell}\left(\psi-k\varphi\right)(\overline{\tau_{1}\dots \tau_{\ell}})}.
\]
Together with \eqref{eq:rn-bdd}, this implies
%\[
%q_N \asymp e^{-NS_{\ell}\left(\varphi\right)(\overline{k_{1}\dots k_{\ell}})}.
%\]
%It follows that 
\[
e^{mS_{\ell}\left(\psi-k\varphi\right)(\overline{\tau_{1}\dots \tau_{\ell}})} \asymp q_m^{k-1} \asymp e^{-m(k-1)S_{\ell}\left(\varphi\right)(\overline{\tau_{1}\dots \tau_{\ell}})},
\]
and hence, 
\[
S_{\ell}\left(\psi-k\varphi\right)(\overline{\tau_{1}\dots \tau_{\ell}})=-(k-1)S_{\ell}\left(\varphi\right)(\overline{\tau_{1}\dots \tau_{\ell}}),
\]
i.e. 
\[
S_{\ell}\left(\psi-\varphi\right)(\overline{\tau_{1}\dots \tau_{\ell}})=0.
\]

We have thus shown that if $\tau_{1}\dots \tau_{\ell}$ contains at least two different letters, then 
\[
S_{\ell}\left(\psi-k\varphi\right)(\overline{\tau_{1}\dots \tau_{\ell}})\neq0\implies S_{\ell}\left(\psi-\varphi\right)(\overline{\tau_{1}\dots \tau_{\ell}})=0.
\]
Therefore, 
\begin{equation}
\label{eq:1or3}
\frac{S_{\ell}\psi\left(\overline{\tau_{1}\dots \tau_{\ell}}\right)}{S_{\ell}\varphi\left(\overline{\tau_{1}\dots \tau_{\ell}}\right)}\in\left\{ 1,k\right\} .
\end{equation}
It remains to consider the fixed points $\overline{\tau_1}$:
\[
\frac{S_{1}\psi\left(\overline{\tau_1}\right)}{S_{1}\varphi\left(\overline{\tau_1}\right)}=\frac{\psi\left(\overline{\tau_1}\right)}{\varphi\left(\overline{\tau_1}\right)}.
\]
Hence, the set
\[
A:=\left\{\frac{S_{\ell}\psi\left(\overline{\tau_{1}\dots \tau_{\ell}}\right)}{S_{\ell}\varphi\left(\overline{\tau_{1}\dots \tau_{\ell}}\right)},\quad \ell \ge 1,\quad \tau_1\dots \tau_\ell \in I \right\}
\]
contains only finitely many values, at least one of which is 1 or $k$ in view of \eqref{eq:1or3}. Now, Lemma \ref{lem:cohom-1} implies
that $\psi$ is cohomologous to $(\dim_{H}\Lambda)\varphi$ and  $\dim_{H}\Lambda=1$.
Hence, $\varphi$ is cohomologous to $\psi$, which is a contradiction. 
\end{proof}

\end{document}
